\exercisesection{1}

¶ 1. (a) Show that a group G cannot be the union of two subgroups distinct from G. Show that, for every set I with at least two elements, the commutative group \(G = \mathbf{F}_{2}^{(I)}\) is the union of three subgroups distinct from G.

\begin{enumerate}
\def\labelenumi{(\alph{enumi})}
\setcounter{enumi}{1}
\item
  Let \((\mathbf{H}_i)_{i\in I}\) be a finite family of subgroups of a group \(G\) such that each of the \(\mathbf{H}_i\) is a subgroup of \(G\) of infinite index. Show that \(G\) cannot be the union of a finite number of left cosets of the \(\mathbf{H}_i\) . (Argue by induction on the number of elements in \(I\) ; if there exist two distinct indices \(i,j\) such that the index \((\mathbf{H}_i:(\mathbf{H}_i\cap \mathbf{H}_j))\) is finite, \(\mathbf{H}_i\) may be suppressed; if on the other hand \(\mathbf{H}_i\cap \mathbf{H}_j\) is of infinite index in \(\mathbf{H}_i\) for every ordered pair \((i,j)\) of distinct indices, consider an index \(k\) such that \(\mathbf{H}_k\) is maximal in the set of \(\mathbf{H}_i\) and show that \(\mathbf{H}_k\) is the union of a finite number of left cosets of the \(\mathbf{H}_k\cap \mathbf{H}_i\) where \(i\neq k\) .)
\item
  Give an example of a commutative ring A and four ideals \(a, b_{1}, b_{2}, b_{3}\) of A such that \(a \notin b_{i} (i = 1, 2, 3)\) , but \(a = \bigcup_{i} b_{i}\) (use (a)).
\end{enumerate}

\begin{enumerate}
\def\labelenumi{\arabic{enumi}.}
\setcounter{enumi}{1}
\item
  Let A be a ring which is not necessarily commutative and \(a, p_{1}, \ldots, p_{n}\) two-sided ideals of A. Suppose that a is contained in the union of the \(p_{i}\) and that all the \(p_{i}\) except at most two are prime ideals (Algebra, Chapter VIII, §8, Exercise 6). Show that a is contained in one of the \(p_{i}\) .
\item
  In the product ring \(\mathbf{A} = \mathbf{R}^{\mathbf{N}}\) , let (for each \(n \in \mathbf{N}\) ) \(\mathfrak{m}_n\) be the maximal ideal consisting of the \(f\colon \mathbf{N} \to \mathbf{R}\) such that \(f(n) = 0\) . Show that \(\mathfrak{a} = \mathbf{R}^{(\mathbf{N})}\) is an ideal of \(\mathbf{A}\) contained in the union of the \(\mathfrak{m}_n\) , but in none of these maximal ideals.
\item
  In a ring A let p be a prime ideal and a an element such that \(p \subset Aa\) but \(a \notin p\) . Show that p = ap.
\item
  \begin{enumerate}
  \def\labelenumii{(\alph{enumii})}
  \tightlist
  \item
    Let A be a Noetherian ring and \(\mathfrak{a}\) an ideal of A. Show that there exists a finite number of prime ideals \(\mathfrak{p}_i\) ( \(1 \leqslant i \leqslant r\) ) such that \(\mathfrak{p}_1\mathfrak{p}_2 \ldots \mathfrak{p}_r \subset \mathfrak{a}\) . (Argue by reductio ad absurdum considering among the ideals containing \(\mathfrak{a}\) , distinct from A and containing no finite product of prime ideals, a maximal element;
  \end{enumerate}
\end{enumerate}

observe on the other hand that, if the ideal b is not prime, there exist two ideals c, b containing b, distinct from b and satisfying cd ⊂ b.) (Cf. Chapter IV.)

If A is an integral domain and \(a \neq 0\) , we may assume that the \(p_{i} \neq 0\) .

\begin{enumerate}
\def\labelenumi{(\alph{enumi})}
\setcounter{enumi}{1}
\tightlist
\item
  Give an example of a non-Noetherian ring A for which the result of (a) is false. (Take for example A to be the ring of continuous real-valued functions on {[}0, 1{]}.)
\end{enumerate}

¶ 6. (a) Let A be a ring and a, b two ideals of A such that b is finitely generated; show that, if the quotient rings A/a and A/b are Noetherian, so is A/ab (observe that b/ab is a finitely generated (A/a)-module).

\begin{enumerate}
\def\labelenumi{(\alph{enumi})}
\setcounter{enumi}{1}
\tightlist
\item
  Show that, if a ring A is such that every prime ideal of A is finitely generated, A is Noetherian. (Argue by reductio ad absurdum by showing that in the set of ideals of A which are not finitely generated there would be a maximal element c which is not prime by hypothesis; there would therefore be two ideals \(a \supset c\) , \(b \supset c\) distinct from c and satisfying \(ab \subset c\) ; then use (a).)
\end{enumerate}

\begin{enumerate}
\def\labelenumi{\arabic{enumi}.}
\setcounter{enumi}{6}
\tightlist
\item
  Let \(\mathbf{A} = \prod_{i=1}^{n} \mathbf{A}_i\) be the product of a finite family of rings, the \(\mathbf{A}_i\) being canonically identified with ideals of \(\mathbf{A}\) . Let \(\mathbf{B}\) be a subring of \(\mathbf{A}\) such that \(\operatorname{pr}_i \mathbf{B} = \mathbf{A}_i\) for \(1 \leqslant i \leqslant n\) .
\end{enumerate}

\begin{enumerate}
\def\labelenumi{(\alph{enumi})}
\item
  Show that, if the \(\mathbf{A}_i\) are Noetherian (resp. Artinian), B is Noetherian (resp. Artinian) (cf.~Algebra, Chapter VIII, § 2, Exercise 12).
\item
  Let \(n_{i}\) be the ideal of B which is the kernel of the restriction of \(pr_{i}\) to B. Show that every prime ideal p of B is contained in one of the \(n_{i}\) (use Proposition 1); deduce that, for this index i, \(pr_{i}\) \(p \neq A_{i}\) . Show that, if each of the \(A_{i}\) contains only a finite number \(k_{i}\) of maximal ideals, B contains at most \(\sum_{i=1}^{n} k_{i}\) maximal ideals.
\item
  Let \(\Re(A)\) , \(\Re(B)\) be the Jacobson radicals of \(A\) and \(B\) respectively. Show that \(\Re(B) = B \cap \Re(A)\) . If each of the \(A_i\) contains only a finite number of maximal ideals, show that, for every integer \(k > 1\) , \(\mathrm{pr}_i((\Re(B))^k) = (\Re(A_i))^k\) for all \(i\) and \((\Re(B))^k = (\Re(A))^k \cap B\) (write \(\Re(B)\) as a product of distinct maximal ideals and note that, if \(m\) is a maximal ideal of \(B\) such that \(\mathrm{pr}_i(m) = m_i\) is a maximal ideal of \(A_i\) , then \(\mathrm{pr}_i(m^k) = m_i^k\) ).
\end{enumerate}

\begin{enumerate}
\def\labelenumi{\arabic{enumi}.}
\setcounter{enumi}{7}
\tightlist
\item
  \begin{enumerate}
  \def\labelenumii{(\alph{enumii})}
  \tightlist
  \item
    Let \(\mathfrak{a}, \mathfrak{b}\) be two relatively prime ideals of \(\mathfrak{a}\) ring \(\mathbf{A}\) . Show that \(\mathfrak{a}: \mathfrak{b} = \mathfrak{a}\) , \(\mathfrak{b}: \mathfrak{a} = \mathfrak{b}\) . If \(\mathfrak{c}\) is an ideal of \(\mathbf{A}\) such that \(\mathfrak{ac} \subset \mathfrak{b}\) , then \(\mathfrak{c} \subset \mathfrak{b}\) .
  \end{enumerate}
\end{enumerate}

\begin{enumerate}
\def\labelenumi{(\alph{enumi})}
\setcounter{enumi}{1}
\item
  Let p, q be two prime ideals neither of which is contained in the other; then p: q = p and q: p = q. Give an example of two principal prime ideals p, q in the polynomial ring A = K{[}X, Y{]} (where K is a field), neither of which is contained in the other and which are not relatively prime.
\item
  Let \(a\) be an element which is not a divisor of 0 in A. Show that, if the principal ideal \(p = Aa\) is prime, the relation \(p = bc\) for two ideals \(b\) , \(c\) of A implies \(b = A\) or \(c = A\) .
\end{enumerate}

\begin{enumerate}
\def\labelenumi{\arabic{enumi}.}
\setcounter{enumi}{8}
\tightlist
\item
  \begin{enumerate}
  \def\labelenumii{(\alph{enumii})}
  \tightlist
  \item
    Let \((\mathfrak{a}_i)_{1 \leqslant i \leqslant n}\) be a finite family of ideals of a ring A which are relatively prime in pairs and such that no \(\mathfrak{a}_i\) is of the form \(\mathfrak{bc} = \mathfrak{b} \cap \mathfrak{c}\) , where \(\mathfrak{b}\) and \(\mathfrak{c}\) are two relatively prime ideals. If \((\mathfrak{a}_j')_{1 \leqslant j \leqslant m}\) is another family of ideals of A which are relatively prime in pairs, such that \(\mathfrak{a}_1\mathfrak{a}_2 \ldots \mathfrak{a}_n = \mathfrak{a}_1'\mathfrak{a}_2' \ldots \mathfrak{a}_m'\) , show that \(m \leqslant n\) ; if \(m = n\) , there exists a permutation \(\pi\) on \([1, n]\) such that \(a_i' = a_{\pi(i)}\) for \(1 \leqslant i \leqslant n\) (use Proposition 5 and Algebra, Chapter VIII, § 1, Exercise 1 (d)).
  \end{enumerate}
\end{enumerate}

\begin{enumerate}
\def\labelenumi{(\alph{enumi})}
\setcounter{enumi}{1}
\tightlist
\item
  Let A be a Noetherian ring. Show that every ideal \(\mathfrak{a}\) of A is equal to the product of a finite number of ideals which are relatively prime in pairs, none of which is the product of two relatively prime ideals. (Argue by reductio ad absurdum considering a maximal element of the set of ideals not having this property.)
\end{enumerate}

\begin{enumerate}
\def\labelenumi{\arabic{enumi}.}
\setcounter{enumi}{9}
\tightlist
\item
  Give an example of a ring A and an infinite family of distinct maximal ideals \((\mathfrak{m}_n)_{n\in \mathbf{N}}\) of A such that \(\bigcap_{n}\mathfrak{m}_{n} = 0\) but the canonical mapping
\end{enumerate}

\[
\mathrm{A} \rightarrow \prod_ {n} (\mathrm{A} / \mathrm{m} _ {n})
\]

is not surjective.

\begin{enumerate}
\def\labelenumi{\arabic{enumi}.}
\setcounter{enumi}{10}
\tightlist
\item
  Let A be a ring which is not necessarily commutative and let \(a, b\) be two two-sided ideals of A such that \(a + b = A\) . Show that \(a \cap b = ab + ba\) . Give an example where \(a \cap b \neq ab\) (consider the ring of lower triangular matrices of order 2 over a field).
\end{enumerate}

